\subsection{Candidate space}

The confirmatory bank evaluated every resolved candidate in each profile ---
\ConfirmResourceConstrainedCandidates{}, \ConfirmIntermediateCandidates{},
and \ConfirmHighCapacityCandidates{} in the resource-constrained,
intermediate-capacity, and high-capacity profiles --- on
\ConfirmResourceConstrainedScenarios{}, \ConfirmIntermediateScenarios{}, and
\ConfirmHighCapacityScenarios{} fresh paired scenarios each, for
\ConfirmTotalExecutions{} executions in total.

The candidate counts fall sharply with profile strength, and that is the
first substantive consequence of separating capacity from controls. A
high-capacity environment already owns least-privilege segmentation and
isolated backups, so those are not decisions it can make. What remains is
patching, detection, isolation, and identity controls:
\ConfirmHighCapacityCandidates{} distinct configurations rather than the
\EnumeratedPortfolios{} the enumeration nominally contains. Under the
previous version's precedence rule the same profile appeared to face far
more choices, many of them cheaper than doing nothing, because a candidate
could sell back the segmentation and backup architecture the profile already
had.

Exactly one candidate per profile costs nothing: the profile's own untouched
posture. Its mean weighted service-hours lost are
\ConfirmResourceConstrainedFreeBaselineMeanLoss{},
\ConfirmIntermediateFreeBaselineMeanLoss{}, and
\ConfirmHighCapacityFreeBaselineMeanLoss{} respectively, and these are the
reference points against which every purchase should be read.

\subsection{Full-space versus finalist-only frontiers}

Over the full space, \ConfirmFullSpaceFrontier{} candidates are
non-dominated: \ConfirmResourceConstrainedNonDominated{},
\ConfirmIntermediateNonDominated{}, and \ConfirmHighCapacityNonDominated{} by
profile. Restricting attention to the \ConfirmFinalistLabelled{} candidates
that the frozen discovery rule labels as finalists, and recomputing dominance
within that subset alone, gives a frontier of \ConfirmFinalistOnlyFrontier{}.

The two views disagree, and how they disagree has changed since the
model was repaired (Figure~\ref{fig:fullspace}).

\textbf{No candidate is efficient only in the restricted view.} The
finalist-only frontier contains \ConfirmFinalistOnlyArtifacts{} candidates
that the full space dominates. An earlier version of this study, run before
the mechanism repairs and without parameter uncertainty, found three such
artifacts; with wider objectives and more scenarios per candidate, that
particular distortion disappears. We report the zero rather than the earlier
three, and note that a design cannot know in advance which it will get.

\textbf{\ConfirmNonFinalistsOnFrontier{} non-dominated candidates were never
labelled as finalists at all.} This is where the restricted design actually
costs something, and it costs more than before: a study evaluating only the
finalist set could not have shown those candidates to a reader under any
circumstances, and could not have reported that they existed. Of the
\ConfirmFinalistLabelled{} labelled finalists, \ConfirmFinalistsSurviving{}
remain non-dominated over the full space and \ConfirmFinalistsDisplaced{} do
not.

The asymmetry is the useful part. A restricted comparison rarely promotes a
candidate that does not deserve it, because the candidates it evaluates were
selected for being good. What it reliably does is hide the ones it never
looked at, and no amount of care inside the restricted set can detect that,
because the evidence needed is exactly the evidence the design discarded.

\begin{figure*}[t]
\centering
\includegraphics[width=\textwidth]{full_space_vs_finalist.pdf}
\caption{Candidate counts at each stage. The gap between the third and fourth
bars is an effect of the candidate set on the reported frontier, not an
effect of the data. Panel subtitles give the number of candidates that appear
efficient within the finalist subset but are dominated once every resolved
candidate is present.}
\label{fig:fullspace}
\end{figure*}

\subsection{Where the decision is live, and where it is not}

The three profiles do not merely differ in degree; they pose different kinds
of decision (Figure~\ref{fig:frontier}).

\textbf{High-capacity: very little worth buying.} All
\ConfirmHighCapacityCandidates{} resolved candidates fall between
\ConfirmHighCapacityMeanLossMin{} and \ConfirmHighCapacityMeanLossMax{}
weighted service-hours. The entire purchasable improvement in this profile is
\ConfirmHighCapacityMaxPurchasableImprovement{} weighted service-hours, and
the best candidate within the \ConfirmHighCapacityBudget-point budget spends
\ConfirmHighCapacityBestAffordableCost{} points to reach
\ConfirmHighCapacityBestAffordableMeanLoss{} hours from a free baseline of
\ConfirmHighCapacityFreeBaselineMeanLoss{}. Against a Monte Carlo standard
error of \ConfirmHighCapacityWorstMeanLossStandardError{} hours, and with
\ConfirmHighCapacityIndistinguishablePairs{} of
\ConfirmHighCapacityCandidatePairs{} candidate pairs
(\ConfirmHighCapacityIndistinguishableShare\%) separated by less than one
combined standard error on every stochastic objective, the ordering within
this profile is weakly determined and should be read that way.

The sensitivity analysis sharpens this into something more interesting than
"the differences are small". In this profile the parameter that explains most
of the variance in mean disruption is not the spread rate or patch
effectiveness but \emph{the false-positive isolation rate}, at
\SensitivityHighCapacityDominantEffect\% of variance. When defenses are
strong enough that the attack rarely spreads far, the dominant modeled source
of lost service-hours is the defense removing healthy systems from service.
That is a conditional model result and not a claim about real hospitals, but
it is the kind of trade-off a six-objective frontier exists to surface, and a
study reporting only mean disruption would have missed it entirely.

\textbf{Intermediate-capacity: the decision that matters.} The free baseline
loses \ConfirmIntermediateFreeBaselineMeanLoss{} hours and the best candidate
within the \ConfirmIntermediateBudget-point budget loses
\ConfirmIntermediateBestAffordableMeanLoss{} at a cost of
\ConfirmIntermediateBestAffordableCost{} points --- a modeled improvement of
more than an order of magnitude, with sustained-outage probability at
\ConfirmIntermediateBestAffordableOutage\%. This is the profile in which the
six objectives genuinely trade against one another and in which a preference
statement changes the answer.

\textbf{Resource-constrained: the budget does not reach.} The free baseline
loses \ConfirmResourceConstrainedFreeBaselineMeanLoss{} hours. The best
candidate anywhere in the space reaches
\ConfirmResourceConstrainedMeanLossMin{}, but the best candidate that fits
the declared \ConfirmResourceConstrainedBudget-point budget reaches only
\ConfirmResourceConstrainedBestAffordableMeanLoss{}, and its sustained-outage
probability is \ConfirmResourceConstrainedBestAffordableOutage\%. Within
budget that endpoint is close to immovable: the affordable options change
mean disruption substantially while leaving multi-service clinical outage
nearly certain. The honest statement is not that some affordable portfolio is
adequate. It is that in this modeled environment the declared budget does not
buy the outcome the endpoint measures.

\begin{figure*}[t]
\centering
\includegraphics[width=\textwidth]{confirmatory_frontier.pdf}
\caption{Cost against mean weighted service-hours lost, with Monte Carlo
standard errors. Coloured markers are candidates non-dominated across all
six objectives; they are deliberately not joined, because a point efficient
in six dimensions can sit above another in this two-dimensional projection
and connecting them would draw a trade-off curve the data do not contain.
The thin step line is the cost/disruption lower envelope: the best mean loss
reachable at or below each cost. Note the vertical scales --- the
high-capacity panel spans under one service-hour.}
\label{fig:frontier}
\end{figure*}

\subsection{Uncertainty and resolution}

Worst-case Monte Carlo standard error on mean disruption is
\ConfirmResourceConstrainedWorstMeanLossStandardError{},
\ConfirmIntermediateWorstMeanLossStandardError{}, and
\ConfirmHighCapacityWorstMeanLossStandardError{} weighted service-hours in
the three profiles.

The share of candidate pairs the scenario bank cannot reliably order is
\ConfirmResourceConstrainedIndistinguishableShare\%,
\ConfirmIntermediateIndistinguishableShare\%, and
\ConfirmHighCapacityIndistinguishableShare\%. The high-capacity figure is
the one to read carefully: at the thirty-scenario pilot resolution that
profile's candidates were indistinguishable to the point of tying exactly,
and at \ConfirmHighCapacityScenarios{} scenarios they separate into sixteen
distinct values while a sixth of pairs still cannot be ordered. Both facts
are about resolution, not about the profile having no decision in it.

Bootstrap Pareto inclusion is at least \StabilityThreshold{} for
\ConfirmResourceConstrainedStableAtEighty{},
\ConfirmIntermediateStableAtEighty{}, and
\ConfirmHighCapacityStableAtEighty{} candidates respectively.

\subsection{The sustained-outage endpoint at every $k$}

Table~\ref{tab:kladder} reports the endpoint at every clinical service count.
The values move by roughly three percentage points across the whole range of
$k$, for the reason given in the methods: within this model the clinical
estate is close to all-or-nothing (Figure~\ref{fig:bimodal}). A reader who
disagrees with the primary $k = \PrimaryK$ can substitute any other and reach
substantially the same conclusion. That is a statement about the model's
behaviour, not a defence of the choice.

\begin{table}[t]
\centering
\caption{Mean sustained-outage probability across candidates, at every
clinical service count $k$. The primary analysis uses $k = \PrimaryK$.}
\label{tab:kladder}
\small
\begin{tabular}{@{}lrrrr@{}}
\toprule
Profile & $k=1$ & $k=2$ & $k=3$ & $k=4$ \\
\midrule
Resource-constrained & \ConfirmResourceConstrainedOutageKOneMean\% & \ConfirmResourceConstrainedOutageKTwoMean\% & \ConfirmResourceConstrainedOutageKThreeMean\% & \ConfirmResourceConstrainedOutageKFourMean\% \\
Intermediate & \ConfirmIntermediateOutageKOneMean\% & \ConfirmIntermediateOutageKTwoMean\% & \ConfirmIntermediateOutageKThreeMean\% & \ConfirmIntermediateOutageKFourMean\% \\
High-capacity & \ConfirmHighCapacityOutageKOneMean\% & \ConfirmHighCapacityOutageKTwoMean\% & \ConfirmHighCapacityOutageKThreeMean\% & \ConfirmHighCapacityOutageKFourMean\% \\
\bottomrule
\end{tabular}
\end{table}

\begin{figure}[t]
\centering
\includegraphics[width=\columnwidth]{clinical_outage_bimodality.pdf}
\caption{Number of clinical services sustaining a qualifying outage across
the discovery bank. The near-binary shape is why the choice of $k$ is not
load-bearing --- and it is an artifact of binary per-step service
availability, not a property of real incidents.}
\label{fig:bimodal}
\end{figure}

\subsection{External validation}

Of \ValidationBenchmarks{} frozen external benchmarks, \ValidationScored{}
could be scored numerically. \ValidationNotAddressable{} are not addressable
because the model has no output of that construct at any parameter setting,
and \ValidationNotScored{} cannot be scored because the model's quantity and
the target's quantity are not the same thing. \ValidationFailed{} benchmarks
fail. None passes.

That ratio is the headline validation result and it bounds every conclusion
above. Three failures are worth stating individually.

\textbf{The restoration engine is anchored to neither reference.} Against the
one directly observed large-scale technical hospital outage the model is far
too slow: its median clinical outage is
\ValidationOutageModelMedianHours{} modeled hours against an observed median
of \ValidationOutageTargetMedianHours{}, and
\ValidationOutageModelSixHourShare\% of disrupted trials recover within six
hours against an observed \ValidationOutageTargetSixHourShare\%
\citep{tully2025}. Against the ransomware incident record the same engine is
far too fast, completing inside a \HorizonHours-hour horizon where the public
duration distribution has most of its mass beyond a week
\citep{neprash2022}. A model simultaneously too slow for a technical outage
and too fast for a ransomware incident is anchored to neither, and its
non-recovery objective should be read as a within-model ordering device
rather than a duration estimate.

\textbf{Multi-facility breadth fails outright.} The simulator runs one
facility, so it assigns probability zero to roughly a quarter of publicly
reported hospital-linked events --- those that affected more than one
Medicare-certified hospital. This is a total discrepancy on a directly
comparable target.

\textbf{Regional spillover cannot be addressed at all.} The model has no
representation of a neighboring hospital, patient transfer, or diversion, so
the documented burden at adjacent emergency departments
\citep{dameff2023,abouk2024} has no counterpart output to compare against.

The one benchmark the model is not obviously wrong on is the share of
incidents contained before clinical service loss, where the model's spread
across profiles brackets the surveyed figure. We report that as consistency
rather than a pass: a self-reported vendor survey spanning many countries and
organization sizes is not a U.S. hospital denominator, and bracketing a point
estimate with a range that wide is weak evidence.

\begin{figure}[t]
\centering
\includegraphics[width=\columnwidth]{external_validation_scorecard.pdf}
\caption{Verdicts against the frozen benchmark registry. Targets and criteria
were fixed before any comparison was run, and failures are reported rather
than corrected.}
\label{fig:validation}
\end{figure}

\subsection{Sensitivity to the unidentified coefficients}

Because every trial records the parameter vector its scenario drew, a
variance decomposition costs no additional simulation: the confirmatory bank
is already a Monte Carlo sample over the declared ranges.
\SampledParameters{} coefficients were sampled.

The results separate two questions that are easy to conflate. In the
resource-constrained and intermediate profiles the dominant parameter for
mean disruption is the base spread rate, explaining
\SensitivityResourceConstrainedDominantEffect\% and
\SensitivityIntermediateDominantEffect\% of variance respectively; in the
high-capacity profile it is the false-positive isolation rate, at
\SensitivityHighCapacityDominantEffect\%.

Those are effects on the \emph{level} of an objective. The effect on the
\emph{decision} is much smaller. Ranking candidates within the low and high
thirds of each parameter and comparing the two orderings gives Spearman
correlations of \SensitivityResourceConstrainedDominantRankStability{},
\SensitivityIntermediateDominantRankStability{} and
\SensitivityHighCapacityDominantRankStability{} for those dominant
parameters, with a worst case of \SensitivityWorstRankStability{} across all
parameters and objectives.

The reading is therefore specific: within the declared ranges, how much
disruption a portfolio produces is genuinely uncertain, while which
portfolios are preferable is largely stable. That is a weaker and more
defensible claim than either "the results are robust" or "the results are
uncertain", and it is the claim the analysis supports. It does not extend
beyond the ranges, which we chose.

\subsection{Structural sensitivity}

Assumption S1 --- that restoration begins only once the estate is contained
--- was varied on the \StructuralCandidatesTested{} candidates carrying the
conclusions above. Structural uncertainty behaves quite differently from
parameter uncertainty, and the contrast is the point of running both.

\textbf{Membership.} \StructuralPreserved{} candidates keep their frontier
membership under both structures; \StructuralOnlyGated{} are non-dominated
only under the primary structure and \StructuralOnlyParallel{} only under the
alternative. Per-profile agreement runs from \StructuralWorstAgreement\% to
\StructuralBestAgreement\%.

\textbf{Magnitude, and its direction.} Removing the containment gate makes
outcomes markedly \emph{worse}, not better. Intermediate-capacity mean loss
across these candidates moves from \StructuralIntermediateGatedMeanLoss{} to
\StructuralIntermediateParallelMeanLoss{} weighted service-hours and
non-recovery from \StructuralIntermediateGatedNonRecovery\% to
\StructuralIntermediateParallelNonRecovery\%; the resource-constrained
profile moves from \StructuralResourceConstrainedGatedMeanLoss{} to
\StructuralResourceConstrainedParallelMeanLoss{}. The high-capacity profile
is unaffected, because it contains the estate almost immediately either way.

The direction is initially counterintuitive and the mechanism is worth
stating, because it is a genuine property of the model rather than an
artifact of the switch. A restored node re-enters the network susceptible.
Restoring while the attack is still propagating therefore returns clean
assets to an adversary who still has reach, and the incident lengthens. S1 is
not a simplifying assumption about scheduling; it encodes the substantive
claim that responders do not restore into an uncontained environment, and the
model is materially more optimistic because of it.

Set against the previous subsection, the comparison is stark. Sampling ten
coefficients across wide declared ranges leaves portfolio orderings largely
intact. Changing one structural assumption moves up to
\StructuralOnlyGated{} candidates off the frontier and nearly doubles modeled
disruption in two profiles. Parameter uncertainty is the kind this study can
quantify; structural uncertainty is the kind that actually threatens its
conclusions, and only one structural assumption of several has been varied.
